\section{Deformations with fixed base and the case \texorpdfstring{$b_2\geq5$}{b2 at least 5}}
\label{def:section}
\label{sec:deformation}

We work in characteristic zero again and write
$f_0:X_0\to B$ for the original projective fibration, with
$e=c_1(f_0^*H)$ and symplectic form $\sigma_0$.
Matsushita's deformation theorem supplies a morphism
defined by a relatively generated line bundle on the entire nearby
$e$-Hodge locus.  We prove that its normal Stein target is analytically
$B\times D$, including when $B$ is singular, by first splitting the
parameter directions locally and then using the rigidity of $B$.
The nearby manifolds need only be compact K\"ahler.

We then follow the inverse image of the fixed closed stratum of
Proposition~\ref{str:minimum} through the family.  When $b_2\geq5$,
this variation forces the Hodge cohomology of $\cS$ to be diagonal,
contradicting Proposition~\ref{eu:euler-zero}.
We first record a restriction calculation that will also control
transported fiber polarizations in Section~\ref{sec:second-system}.

\begin{lemma}\label{def:restriction-kernel}
Let $i:F\hookrightarrow X_0$ be a smooth fiber over the smooth locus of
$f_0$.  Then
\begin{equation}\label{def:kernel-equation}
 \ker\bigl(i^*:H^2(X_0,\C)\longrightarrow H^2(F,\C)\bigr)=e^\perp.
\end{equation}
\end{lemma}

\begin{proof}
The smooth fibers have one common cohomology class, and intersecting
$n$ general hyperplanes of the base gives
$e^n=\deg_H(B)[F]$.  Polarizing the Fujiki relation, and using $q(e)=0$,
gives a positive constant $a$ such that
\begin{equation}\label{def:mixed-fujiki}
 \int_F i^*\xi_1\cdots i^*\xi_n
       =a\prod_{j=1}^n q(e,\xi_j)
       \qquad(\xi_j\in H^2(X_0,\C)).
\end{equation}
Indeed, in every nonzero pairing of the $2n$ entries
$e,\ldots,e,\xi_1,\ldots,\xi_n$, each copy of $e$ must be paired
with a different $\xi_j$.  Positivity of $a$ follows by taking all
$\xi_j$ to be an ample class $h$; in particular $q(e,h)>0$.

The restrictions of $\sigma_0$ and $\bar\sigma_0$ vanish.  If
$\xi\in H^{1,1}(X_0,\R)$ and $q(e,\xi)=0$, Equation~\eqref{def:mixed-fujiki}
gives
\[
 \int_F i^*\xi\,(i^*h)^{n-1}=0,
 \qquad
 \int_F(i^*\xi)^2(i^*h)^{n-2}=0
\]
when $n\geq2$.  The Hodge index theorem on $F$ makes the second
quadratic form negative definite on the primitive real $(1,1)$-classes,
so $i^*\xi=0$.  For $n=1$, the first integral already proves this,
since $H^2(F,\R)$ is one-dimensional.  Thus $e^\perp$ is contained
in the kernel.  Conversely, the first equality in
Equation~\eqref{def:mixed-fujiki}, with one entry $\xi$ and all other
entries $h$, shows that every class whose restriction vanishes is
orthogonal to $e$.  The Hodge decomposition and complexification give
Equation~\eqref{def:kernel-equation}.
\end{proof}

Let $D$ be a sufficiently small polydisc about the original point in the
locus $\operatorname{Def}(X_0,e)$ on which $e$ is of type $(1,1)$.
The local Torelli theorem identifies this smooth germ with the period
hyperplane $q(e,\sigma_t)=0$.  Write
$\pi:\mathcal X\to D$ for the restricted Kuranishi family, and use its
cohomology marking throughout.

\begin{proposition}[Fixed-base deformation]\label{def:fixed-base}
After shrinking $D$, there is a proper holomorphic map
\[
             \mathfrak f:\mathcal X\longrightarrow B\times D
\]
over $D$ whose central fiber is $f_0$.  Every
$f_t:X_t\to B$ has connected fibers, and every irreducible component
of every fiber is an $n$-dimensional Lagrangian subvariety.
Under the chosen marking, $c_1(f_t^*H)=e$ for every $t$.
For some $m>0$, the extension of $f_0^*H^m$ is relatively generated
and defines these fibrations after Stein factorization.
\end{proposition}

\begin{proof}
Choose $m>0$ so that $H^m$ is very ample, and put
$L_0=f_0^*H^m$.  The local universal deformation of the pair
$(X_0,L_0)$ is the restriction of the Kuranishi family to
\[
                   q(c_1(L_0),\sigma_t)=m q(e,\sigma_t)=0,
\]
and carries an extension $\mathcal L$ of $L_0$; this is
\cite[Theorem~1.1]{MatsushitaDeformations2016}.  Thus its restriction
to a small polydisc is exactly the family $\mathcal X\to D$ chosen
above.  Theorem~1.2 and Corollary~1.3 of the same source give a locally
free sheaf $\mathcal E=\pi_*\mathcal L$ and a surjective evaluation
$\pi^*\mathcal E\to\mathcal L$.  Their hypotheses are the original
Lagrangian fibration over the projective $B$ and
$L_0=f_0^*H^m$ with $H^m$ very ample; the nearby fibers may be nonprojective compact
K\"ahler manifolds.  The resulting proper morphism is
\[
                  \Phi:\mathcal X\longrightarrow\PP_D(\mathcal E).
\]
Here $\PP_D(\mathcal E)$ parametrizes one-dimensional quotients.
The cohomology-and-base-change argument in Theorem~1.2 identifies
$\mathcal E\otimes k(0)$ with $H^0(X_0,L_0)$.  Since
$f_{0*}\cO_{X_0}=\cO_B$, projection formula identifies this last
space with $H^0(B,H^m)$.  Hence $\Phi_0$ is $f_0$ followed by the
very ample embedding of $B$ defined by the complete system $|H^m|$.

Let $\mathcal Y$ be the image of $\Phi$, with its reduced analytic
structure, and let
\[
 \mathcal X\xrightarrow{F}\mathcal B\xrightarrow{\nu}
             \mathcal Y\subset\PP_D(\mathcal E)
\]
be its relative Stein factorization.  Write $p:\mathcal B\to D$.
The space $\mathcal B$ is normal because $\mathcal X$ is normal,
and $p$ is proper.  For each $t$, take the normal Stein factorization
of the induced map to its image,
\[
              X_t\xrightarrow{g_t}Z_t
                    \longrightarrow\Phi_t(X_t).
\]
Here $Z_t$ is normal and finite over a projective variety, and $g_t$
has connected fibers.  Corollary~1.3 gives a general Lagrangian fiber.
The symplectic form on the compact K\"ahler manifold $X_t$ is
$d$-closed.  Thus the proper map $g_t$ satisfies
\cite[Theorem~1]{Matsushita2000}, which makes every irreducible
component of every fiber an $n$-dimensional Lagrangian subvariety.
The underlying set of each fiber of $F$ is a connected component
of a fiber of $\Phi_t$, and hence is the underlying set of a fiber
of $g_t$.  In particular $F$ is equidimensional of relative
dimension $n$.  This determines
its fiber dimensions before any identification of a fiber of
$\mathcal B$ with a fiberwise Stein target.

We prove that $p$ is locally an analytic product before making any
base-change assertion for its central fiber.  Fix
$b\in p^{-1}(0)$ and $x\in F^{-1}(b)$.  A general codimension-$n$
smooth slice through $x$ in $X_0$ meets $F^{-1}(b)$ only at $x$
as a germ.  Extend its local defining equations to $\mathcal X$.
The resulting slice $Y$ is smooth over $D$, and its map to the germ
$(\mathcal B,b)$ is finite and dominant.  This follows from the
isolated-fiber criterion for a finite map of analytic germs and
equality of dimensions.  Let $R\subset A$ be the corresponding
finite extension of local rings and let $r$ be its generic degree.
Here $R$ is normal, $A$ is regular, and the parameter functions
$t_1,\ldots,t_s$ are part of a smooth coordinate system on $Y$.

Choose derivations $\delta_i$ of $A$ satisfying
$\delta_i(t_j)=\delta_{ij}$.  Define
\begin{equation}\label{def:trace-derivations}
 \partial_i(a)=\frac1r
       \operatorname{Tr}_{\operatorname{Frac}A/
                             \operatorname{Frac}R}
                    \bigl(\delta_i(a)\bigr),\qquad a\in R.
\end{equation}
Trace takes $A$ into $R$: its values are integral over $R$ and
belong to its fraction field, and $R$ is normal.  Trace is
$R$-linear, so the Leibniz rule shows that $\partial_i$ is a
derivation of $R$, with $\partial_i(t_j)=\delta_{ij}$.
Holomorphic derivations of an analytic space have local flows.
For example, embed the germ in a polydisc, lift the derivation to a
vector field preserving its defining ideal, and restrict the
ambient flow.  Successive flows of the $\partial_i$ give an
isomorphism of germs
\begin{equation}\label{def:base-local-product}
             (\mathcal B,b)\simeq(p^{-1}(0),b)\times(D,0)
\end{equation}
over $D$.  The lifts need not commute: inverse flows in the reverse
order give the inverse product map.

Normality of $\mathcal B$ now implies normality of $p^{-1}(0)$.
The central factorization has connected first fibers and a finite
second map, and its composite is $\Phi_0$, the complete system
$|H^m|$ pulled back by $f_0$.  Uniqueness of Stein factorization
therefore identifies $p^{-1}(0)$ with the original $B$.

It remains to globalize Equation~\eqref{def:base-local-product}.
It gives an exact sequence
\[
 0\longrightarrow T_{\mathcal B/D}
 \longrightarrow T_{\mathcal B}
 \longrightarrow p^*T_D\longrightarrow0.
\]
The relative tangent sheaf is flat over $D$ and restricts to $T_B$,
as is seen in the local products.  By
Corollary~\ref{cor:tangent-rigidity}, $H^1(B,T_B)=0$.
Proper coherent semicontinuity consequently gives
$R^1p_*T_{\mathcal B/D}=0$ after shrinking.  Also
$p_*\cO_{\mathcal B}=\cO_D$, since each fiber of $p$ is the image
of the connected $X_t$ and is therefore connected.
Pushing forward the tangent sequence and using that $D$ is Stein
therefore gives global holomorphic lifts of the coordinate vector
fields of $D$.  Properness supplies one neighborhood on which
their successive flows are defined along the whole central fiber.
These flows give $\mathcal B\simeq B\times D$ over $D$.
The resulting $\mathfrak f$ has the asserted fiber properties.
The line bundle $\mathfrak f^*\operatorname{pr}_B^*H$ on
$\mathcal X$ restricts centrally to $f_0^*H$, so its fiber Chern
classes are the flat class $e$ under the chosen marking.
\end{proof}

\begin{remark}\label{def:nonprojective-warning}
The vanishing $H^1(B,T_B)=0$ comes from the original projective
fibration.  Proposition~\ref{def:fixed-base} produces proper
holomorphic fibrations on every small $e$-deformation, whose sources
may be nonprojective compact K\"ahler manifolds.  Thus the next
argument uses an open set in the full period hyperplane.
Section~\ref{sec:second-system} will select a projective member by
constructing an integral K\"ahler class.
\end{remark}

Assume that a nonquotient germ remains, and retain the closed stratum
$S_0$, its smooth stack $\cS$, and $d=\dim\cS<n$ from
Proposition~\ref{str:minimum}.  The local charts in
Proposition~\ref{loc:maximal-charts} identify the preimage of this
stratum with its smooth orbit on each maximal chart.

\begin{lemma}\label{def:tangent-to-orbit}
On a maximal chart, let $g:Y\to V$ be one of the lifted smooth
source maps, and let $S\subset V$ be the orbit stratum.  At every
$y\in g^{-1}(S)$,
\begin{equation}\label{def:differential-tangency}
                  dg_y(T_yY)\subset T_{g(y)}S.
\end{equation}
The same assertion holds on each parameter fiber of the family in
Proposition~\ref{def:fixed-base}.
\end{lemma}

\begin{proof}
First take a smooth slice $T$ through $y$ finite over $(V,g(y))$.
For a vector $v\in T_yT$, choose a local derivation $\delta$ of
$T$ with that value.  The normalized trace construction in
Equation~\eqref{def:trace-derivations} gives a derivation
$\partial$ of $V$.  Because the chosen local finite branch has
only $y$ over $g(y)$, the trace of a function specializes to its
value at $y$ times the local degree.  Thus
$\partial(g(y))=dg_y(v)$.  Evaluations of all local derivations of
the chart are precisely the tangent space of its orbit, by
Proposition~\ref{str:minimum}.  This proves
Equation~\eqref{def:differential-tangency} for vectors tangent to
the slice.

The tangent spaces of finite smooth slices span $T_yY$.
Indeed, the condition that a complementary linear slice miss the
projectivized fiber tangent cone is a nonempty open condition in
the appropriate Grassmannian; its tangent planes cannot all lie
in a proper linear subspace.  Taking these slices proves the
claim for every vector.  The argument is analytic and applies
unchanged to the compact K\"ahler deformations.  The normalized
source covers used here are smooth: their maps to $\mathcal X$
are finite and unramified outside a set of codimension at least
two, by equidimensionality, and hence are \'etale by purity.
\end{proof}

Resolve the reduced inverse image of $S_0\times D$ in $\mathcal X$
by a resolution functor compatible with smooth maps, using the
complex analytic form of \cite[Theorem~1.1]{BierstoneMilman2008}.
After shrinking around the compact central fiber, finitely many
blowups suffice.  Discard components
whose image in $D$ is proper, and then restrict to a nonempty open
$D^\circ\subset D$ over which the resolution is a smooth proper
family.  Its fibers will be denoted $Q_t$.  They are compact
K\"ahler manifolds: they are projective modifications of analytic
subspaces of the compact K\"ahler fibers $X_t$.
The lifted chart torsors descend to maps
\[
       j_t:Q_t\longrightarrow X_t,
       \qquad h_t:Q_t\longrightarrow\cS,
\]
and some component of $Q_t$ dominates $S_0$.
We may allow $Q_t$ to be disconnected.

\begin{lemma}[Hamiltonian map along the stratum]\label{def:hamiltonian}
Put $\omega_t=j_t^*\sigma_t$.  There is a holomorphic bundle map
\begin{equation}\label{def:hamiltonian-equations}
 \mathcal H_t:h_t^*\Omega^1_{\cS}\longrightarrow T_{Q_t},
 \qquad
 \iota_{\mathcal H_t(\theta)}\omega_t=h_t^*\theta,
 \qquad dh_t\circ\mathcal H_t=0.
\end{equation}
It is independent of local extensions of $\theta$ to ambient base
forms, and the maps are compatible with the family over $D^\circ$.
\end{lemma}

\begin{proof}
On a maximal chart, lift a form on the smooth orbit to an ordinary
holomorphic one-form $\widetilde\theta$ on the ambient base chart.
On its smooth source define the relative Hamiltonian vector field
$v$ by
$\iota_v\sigma_t=g^*\widetilde\theta$.
It is vertical on the dense locus where the fibration is smooth,
because those fibers are Lagrangian, and hence it is vertical
everywhere.  Equivalently, Hamiltonians of two functions pulled
back from the base have identically zero Poisson bracket.
Thus the local flow preserves the reduced inverse image of the
stratum.

If two extensions of $\theta$ are chosen, their difference belongs
to $I\Omega^1_V+dI$, where $I$ is the ideal of the orbit.
Its pullback as a covector on the smooth source is zero at every
point over the orbit, by Lemma~\ref{def:tangent-to-orbit}.
Nondegeneracy of the ambient symplectic form then says that the
difference of the two vector fields is zero on this reduced inverse
image.  Consequently they induce the same derivation there.

Functoriality for the smooth joint flow map and for the product with
its parameter disc lifts the relative flow on the reduced total inverse
image to the total analytic resolution over $D^\circ$.  The lifted
flow fixes the deformation parameter, so differentiation gives relative
vector fields there.  Additivity, $\cO$-linearity, and
compatibility on overlapping charts hold on the dense isomorphism
locus of the total resolution.  Their differences are holomorphic
relative vector fields, so these identities extend over that total
resolution.  This gives a relative bundle map.  The contraction and
verticality identities also hold on the same dense locus and extend
as holomorphic relative tensor identities over the total resolution.
Restricting the map and identities to $Q_t$ gives
Equation~\eqref{def:hamiltonian-equations}.  In particular the image
vectors are pairwise isotropic and annihilate every form pulled back
from $\cS$.
\end{proof}

We recall the Hodge-theoretic meaning of this smooth stack.  Its rational
and complex cohomology are those of $S_0$: finite stabilizers have no
positive-degree cohomology with these coefficients, whether or not
they act effectively \cite[Proposition~36]{Behrend2004}.
Moreover $\cS$ has a K\"ahler orbifold metric.
For example, pull back a Fubini--Study form from a projective embedding
of $S_0$.  In a finite chart its potential has a strict minimum at the
marked point.  The regularized maximum with a sufficiently small
invariant positive quadratic potential, shifted by a small positive
constant, agrees with the old potential near the chart boundary and
is strictly plurisubharmonic near that point.  It remains semipositive
elsewhere.  Finitely many such modifications and their sum give a
positive form everywhere.  The compact K\"ahler orbifold Hodge theorem gives Hodge decomposition
and the K\"ahler identities in this proper \'etale groupoid setting
\cite[Definition~5.23 and Theorems~5.26--5.28]{MilanovSaito}.
The complex orientation and properness also give Poincar\'e duality
\cite[Corollary~25]{Behrend2004}.  We use this ordinary cohomology, rather than an inertia-weighted
Euler characteristic.

\begin{proposition}\label{def:large-betti}
If $b_2(X)\geq5$, the base $B$ has only quotient singularities.
\end{proposition}

\begin{proof}
Suppose a stratum $\cS$ as above exists.  We prove
\begin{equation}\label{def:diagonal-hodge}
                 H^{p,q}(\cS)=0\qquad(p\ne q).
\end{equation}
If $d=0$, its coarse space is a point, already contradicting
Proposition~\ref{eu:euler-zero}; we may assume $d>0$.

Fix $\alpha\in H^{p,q}(\cS)$ with $p>q$, and set
$k=d-p+1$.  On every $Q_t$,
\begin{equation}\label{def:holomorphic-vanishing}
                    \omega_t^k\,h_t^*\alpha=0.
\end{equation}
To verify it, first note that the pullback of $\sigma_t$ by every
smooth space mapping into a fiber of $f_t$ is zero.  On a resolution
of each irreducible fiber component this follows from the
Lagrangian assumption and extension of holomorphic forms.
For a map whose image is contained in a proper subvariety of that
component, resolve a component of its proper base change to that
resolution which dominates the source.  Pullback of forms by a
dominant map is injective, giving the same conclusion.
Thus, at a general point where $h_t$ has constant rank,
$\omega_t$ vanishes on two vertical vectors.  Each of its terms
has at least one horizontal one-form factor.  Wedge multiplication
by $\omega_t^k$ consequently kills the pullback of every
holomorphic $p$-form on $\cS$, since $k+p>d$.  This equality of
holomorphic bundle maps extends from the dense open set and
implies Equation~\eqref{def:holomorphic-vanishing} on Dolbeault
cohomology.

Work over a simply connected small open subset of $D^\circ$.
Smooth proper transport makes the maps $j_t^*$ and $h_t^*$
constant on cohomology.  Accordingly
\[
 v\longmapsto (j_t^*v)^k\,h_t^*\alpha
\]
is one fixed homogeneous polynomial with values in a fixed
cohomology group.  It vanishes on all period vectors over that
open set, by Equation~\eqref{def:holomorphic-vanishing}.
The quadratic form on $e^\perp_\C$ has radical $\C e$ and rank
$b_2-2\geq3$.  Its null cone is therefore irreducible.  Local
Torelli and Proposition~\ref{def:fixed-base} imply that the
periods just used contain an analytic open subset of its
projectivization.  Such an open set is Zariski dense.  Hence the
polynomial vanishes on the whole cone, including the conjugate
period of our chosen fiber.  Suppressing $t$, we have obtained
\begin{equation}\label{def:antiholomorphic-vanishing}
                       \bar\omega^{\,k}h^*\alpha=0.
\end{equation}

We give the contraction step with its coefficients retained.
This avoids replacing a sheaf-valued class by its exterior product
before the contractions have been applied.  If $\alpha\ne0$,
Poincar\'e duality supplies
$\beta\in H^{d-p,d-q}(\cS)$ with $\alpha\beta\ne0$.
Pass to the conjugate complex structures and put
\[
 \mathcal A=h^*\Omega^1_{\overline{\cS}},
 \qquad r=d-q,
 \qquad s=\bar\omega\in H^0(\bar Q,\Omega^2_{\bar Q}).
\]
Write $\rho:\mathcal A\to\Omega^1_{\bar Q}$ for differential
pullback.
In this notation
\[
 h^*\alpha\in H^p(\bar Q,\Omega^q_{\bar Q}),
 \qquad
 \widetilde\beta=h^*\beta
       \in H^{d-p}(\bar Q,\Lambda^r\mathcal A).
\]
The second class is kept in the indicated coefficient bundle;
it has not yet been wedge-mapped to $\Omega^r_{\bar Q}$.
Since $p>q$, we have $r\geq k$.
Apply the exterior coproduct
\[
 \Delta_{k,r-k}:\Lambda^r\mathcal A
       \longrightarrow\Lambda^k\mathcal A
                         \otimes\Lambda^{r-k}\mathcal A
\]
to its coefficients.  Cup Equation~\eqref{def:antiholomorphic-vanishing}
with this class, contract the first $k$ coefficient factors by
$\overline{\mathcal H}$ from
Lemma~\ref{def:hamiltonian}, and wedge the remaining $r-k$
factors as basic forms.  All these are holomorphic sheaf maps on
$\bar Q$, so they induce well-defined operations on its sheaf
cohomology.  Their output belongs to
$H^d(\bar Q,\Omega^d_{\bar Q})$.

Here is the exterior-algebra computation of that output.  If
$a_1,\ldots,a_k$ are local sections of $\mathcal A$ and $\eta$
is a basic $q$-form, the identities
\[
 \iota_{\overline{\mathcal H}(a)}s=\rho(a),
 \qquad
 \iota_{\overline{\mathcal H}(a)}\rho(b)=0
\]
give
\[
 \iota_{\overline{\mathcal H}(a_k)}\cdots
 \iota_{\overline{\mathcal H}(a_1)}(s^k\wedge\eta)
 =(-1)^{k(k-1)/2}k!\,
       \rho(a_1)\wedge\cdots\wedge\rho(a_k)\wedge\eta.
\]
For a decomposable $r$-form the signed shuffle coproduct followed
by wedge multiplication is $\binom rk$ times the original form.
It follows, with only an overall sign depending on conventions,
that the preceding sheaf cup operation is
\begin{equation}\label{def:contraction-output}
                      \pm k!\binom rk\,h^*(\alpha\beta).
\end{equation}
Its input was zero, so Equation~\eqref{def:contraction-output}
forces $h^*(\alpha\beta)=0$.

That is impossible.  Choose a component $Q'$ dominating $S_0$,
a K\"ahler form $\kappa$ on $Q'$, and a K\"ahler orbifold form
$\tau$ on $\cS$.  If $N=\dim Q'$, then
\[
                 \int_{Q'}h^*\tau^d\wedge\kappa^{N-d}>0:
\]
the integrand is nonnegative and positive on the nonempty open
set where $h$ has rank $d$.  Thus the top class of $\cS$ pulls
back nontrivially.  As $\alpha\beta$ is a nonzero top class, it
cannot pull back to zero.  This proves the vanishing for $p>q$;
complex conjugation proves Equation~\eqref{def:diagonal-hodge}.

The Euler number is now
\[
             \chi_{\mathrm{top}}(S_0)
                =\sum_{p=0}^d h^{p,p}(\cS)>0,
\]
because the coarse stratum is nonempty and connected.  This
contradicts Proposition~\ref{eu:euler-zero}.  There can therefore
be no nonquotient germ.
\end{proof}

\begin{remark}\label{def:rank-four-boundary}
The irreducibility just used is the only place in this section where
$b_2\geq5$ is required.  When $b_2=4$, the restricted quadratic form
has rank two, and its null cone is the union of two distinct linear
components.  A local period branch lies in one component and its
conjugate in the other.  The continuation from
Equation~\eqref{def:holomorphic-vanishing} to
Equation~\eqref{def:antiholomorphic-vanishing} then does not follow.
The two rank-four arguments below are necessary for exactly this
reason.
\end{remark}
